\documentclass[11pt,twoside]{book}
\usepackage[T1]{fontenc}
\usepackage[utf8]{inputenc}
\usepackage{microtype}
\usepackage{geometry}
≥ometry{
    paperwidth=6in,
    paperheight=9in,
    margin=1.1in,
    headheight=14pt
}
\usepackage{lettrine}
\usepackage{titlesec}
\usepackage{fancyhdr}
\usepackage{enumitem}
\usepackage{tabularx}
\usepackage{booktabs}
\usepackage{longtable}
\usepackage[most]{tcolorbox}
\tcbuselibrary{breakable}
\tcbset{breakable}
\usepackage{adforn}
\usepackage{hyperref}
\usepackage{setspace}
\usepackage{parskip}
\usepackage{epigraph} % Added for epigraphs
\setlength{\epigraphwidth}{0.7\textwidth}
\renewcommand{\epigraphflush}{flushright}

% Colors
\usepackage[dvipsnames]{xcolor}
\definecolor{royalblue}{RGB}{65,90,140}
\definecolor{agedpaper}{RGB}{250,245,235}

% Typography − A classic Garamond for that legal/manuscript feel
\usepackage{ebgaramond}
\renewcommand{\rmdefault}{ebgaramond}

% Chapter styling − formal, with rules
\titleformat{\chapter}[display]
{\normalfont¢ering\Huge\bfseries\color{royalblue}}
{\chaptertitlename\ þechapter}{20pt}{\Huge}
\titlespacing*{\chapter}{0pt}{0pt}{40pt}

% Section styling
\titleformat{\section}
{\normalfontŁarge\bfseries\scshape}{þesection}{1em}{}
\titleformat{\subsection}
{\normalfontłargeℹtshape}{þesubsection}{1em}{}

% Header/footer
\pagestyle{fancy}
\fancyhf{}
\fancyhead[LE]{þepage \textsc{\nouppercase{≤ftmark}}}
\fancyhead[RO]{\textsc{\nouppercase{\rightmark}} þepage}
\fancyhead[RE,LO]{}
\renewcommand{ℎeadrulewidth}{0.4pt}
\renewcommand{\footrulewidth}{0pt}

% Custom environments
≠wenvironment{sidebar}[1][]
{\begin{tcolorbox}[colback=agedpaper,colframe=royalblue,boxrule=0.5pt,arc=2mm,title={\scshape\small #1},fonttitle=\bfseries]}
{\end{tcolorbox}}

≠wenvironment{quotebox}
{\begin{quote}\smallℹtshape}
{\end{quote}}

% List styling
\setlist[itemize]{noitemsep,topsep=0.5em}
\setlist[enumerate]{noitemsep,topsep=0.5em}

% Hanging indents for legal−style lists (improved)
≠wcommand{≤galitem}[1]{\par\noindent\textbf{#1}\par\noindentℎangindent=2emℎangafter=1}

% Ornament macro
≠wcommand{ørnament}{\par¢ering\adforn{36}\par}

\title{\Huge\scshape The Hedge−Law Realm\\\vspace{0.5em}Łarge\normalfont\textit{A Viterran Expansion for Fate's Edge}}
\author{\textit{Set down by Dusana of the Raven Road, under a borrowed lamp in the White Hall}\\[0.5em]\small\textit{The ink is gall and the ash of a burned writ. The binding is leather from a debtor's cow.}}
\date{}

\begin{document}

\frontmatter

\maketitle

þispagestyle{empty}
\begin{center}
\vspace*{2em}
\textit{The ink is gall and the ash of a burned writ.}\\
\textit{The binding is leather from a debtor's cow.}
\vspace{2em}
\end{center}

\tableofcontents

\mainmatter

\chapter[Preface: The Queen's Peace]{Preface: The Queen's Peace}

≤ttrine[lines=3,findent=2pt]{T}{he law is not a sword.} It is a scale. A sword can be stolen; a scale only tilts.

\begin{quotebox}
−−− Queen Veralyn Everblood, from her first address to the Parlement
\end{quotebox}

Viterra is not a land of crumbling ruins or haunted moors. It is a land of hedgerows and courthouses, of ancient customs written in margins, of disputes settled by gavel and precedent rather than by blade. The Queen's Peace is not a promise−−−it is a contract, renewed each morning by every magistrate, every justicar, every clerk who opens a ledger.

After the succession wars that bled the eastern marches for a generation, Queen Veralyn Everblood sits on a throne still warm from her rivals' blood. She is young, unmarried, and surrounded by counselors who mistake caution for loyalty. Her reign is defined by three pressures:

\begin{itemize}
    ℹtem The need to rebuild law where war left only custom.
    ℹtem The presence of turncloak nobles who switched sides multiple times during the wars and now demand pardon.
    ℹtem The quiet power of the Justicars−−−magistrates who can travel anywhere, hear any case, and speak with the Queen's own voice.
\end{itemize}

This expansion provides the tools to play in Viterra: legal duels, oath−courts, hedge−law, and the politics of a realm that has learned, at great cost, that justice is the only thing that holds a kingdom together when armies have gone home.

ørnament

\chapter{The Pillars of Viterran Law}

Viterran law is not written in a single codex. It lives in four sources, each with its own weight.

\section{The Queen's Word}

The monarch may issue decrees, but a decree that contradicts the \textit{Codex of Precedents} is null unless confirmed by the Parlement. Queen Veralyn has issued only seven decrees in five years; her predecessor issued forty in a single session.

A royal decree carries the weight of the crown's authority. To defy it is treason. But the Queen herself may be petitioned to set aside a decree that causes manifest harm−−−a check on absolute power that has not been used in living memory.

\section{The Codex of Precedents}

A massive book, kept in the White Hall under lock and witness, recording every judgment of the Queen's Justicars for four centuries. When a new case arises, the first question is always: \textit{"What does the Codex say?"} If no precedent exists, the Justicar creates one−−−and the next Justicar must follow it.

The Codex is not infallible. Pages have been lost to fire, water, and deliberate excision. Some precedents contradict one another; older clerks whisper that the Codex sometimes \textit{changes} overnight, a new clause appearing in a margin as if written by a ghost.

\begin{sidebar}[Strings (Viterran Law)]
\begin{itemize}
    ℹtem \textbf{Precedent Citation} −− Reroll a failed Insight roll related to law or custom.
    ℹtem \textbf{Codex Page} −− A copied folio; once per session, reduce DV by 1 on a legal argument.
\end{itemize}
\end{sidebar}

\section{The Hedge−Law}

The unwritten customs of villages and rural manors. A hedge−court has no gavel, only a stool and a witness. Its judgments are binding only on those who consent−−−but refusal to consent is itself a mark against one's reputation. Hedge−law covers boundary stones, water rights, grazing rotations, and the proper distance between beehives.

Hedge−courts are held under ancient oaks or at crossroads. The judge is often the oldest person present, or the one with the most even temper. Verdicts are pragmatic: a stolen pig is replaced; a broken fence is mended by the one who broke it; a blood−feud is settled by a marriage.

\begin{sidebar}[Strings (Hedge−Law)]
\begin{itemize}
    ℹtem \textbf{Boundary Stone} −− Declare a disputed line once; the hedge−court will honor your claim for a season.
    ℹtem \textbf{Hedge Writ} −− Safe passage through rural lands; no tolls, no questions.
\end{itemize}
\end{sidebar}

\section{The Right of Acknowledgment}

The newest pillar, added by Queen Veralyn's own hand after the trial of Baron Aldric the Turncloak. It permits any child who was not ransomed, not visited, or not named in a will to demand that a parent speak one truth aloud under witness.

The truth need not be kind. It need only be honest.

Once spoken, the debt is considered paid in full. The parent cannot be compelled again. The child may walk away−−−or stay. The law does not care.

Mechanically, invoking the Right of Acknowledgment is a social action (Presence + Sway or Command) against the parent's Resolve. On a hit, the parent speaks a true answer to one question. On a critical, the answer must be complete and without evasion. The parent may refuse, but refusal is an admission of guilt in the court of public reputation−−−they become \textit{anima non grata}, and no lawful burial ground will accept their body.

ørnament

\chapter{Factions of Viterra}

\section{The Queen's Justicars}

A small order of traveling magistrates, each bearing a silver gavel and a sealed writ. A Justicar may hear any case, anywhere, and her judgment is final unless overturned by the Queen herself. There are only twelve Justicars in all of Viterra. They are feared, respected, and occasionally murdered.

\begin{sidebar}[Strings]
\begin{itemize}
    ℹtem \textbf{Gavel token} (force a parley under witness)
    ℹtem \textbf{Justicar's seal} (DV −1 for legal negotiations)
    ℹtem \textbf{Precedent citation} (reroll a failed Insight roll related to law)
\end{itemize}
\end{sidebar}

\section{The Hedge−Courts}

Local assemblies of elders, millers, and retired soldiers who settle disputes without reference to the Codex. Their judgments are pragmatic, swift, and often brutal. A hedge−court can outlaw a person from a village−−−a sentence of slow starvation.

\begin{sidebar}[Strings]
\begin{itemize}
    ℹtem \textbf{Boundary stone} (declare a disputed line once)
    ℹtem \textbf{Hedge writ} (safe passage through rural lands)
    ℹtem \textbf{Goose feather} (one use to call a moot)
\end{itemize}
\end{sidebar}

\section{The Turncloak Nobility}

Nobles who changed sides during the succession wars. Some did so out of opportunity; others out of fear. Queen Veralyn has pardoned most but forgotten none. Their estates are intact, but their honor is stained. Many seek redemption through service−−−or revenge through conspiracy.

\begin{sidebar}[Strings]
\begin{itemize}
    ℹtem \textbf{Pardon letter} (reduces Exposure by 2 in Viterra)
    ℹtem \textbf{Oath of fealty} (once/arc, call on a turncloak for aid, but mark a Debt to them)
    ℹtem \textbf{Blood−silver receipt} (proof of payment for a past betrayal)
\end{itemize}
\end{sidebar}

\section{The Parlement of Viterra}

An assembly of nobles, bishops, and guildmasters that meets twice a year to advise the Queen and approve taxes. The Parlement has no power to make law, but it can refuse to fund the Queen's projects. Veralyn has learned to charm, bribe, and threaten them by turns.

\begin{sidebar}[Strings]
\begin{itemize}
    ℹtem \textbf{Proxy vote} (once/session, sway a Parlement decision)
    ℹtem \textbf{Tax receipt} (reduces Supply cost by 1 in Viterran markets)
    ℹtem \textbf{Censure motion} (public shaming of an enemy)
\end{itemize}
\end{sidebar}

ørnament

\chapter{New Talents: The Justicar's Path}

For characters who walk the road of law, oath, and judgment. These talents are available to any class, but they resonate most strongly with Runekeepers of Mykkiel or The Witness.

\section{Foundation Talents (2−−4 XP)}

≤galitem{Cite Precedent (2 XP)} Once per session, you may declare that a past judgment covers the current situation. Roll Wits + Lore (DV 3). On a hit, gain +1 die on your next legal or social roll. On a miss, the precedent is contested−−−GM gains 1 SB.

≤galitem{Gavel's Authority (4 XP)} You carry a small gavel (or a symbolic equivalent). When you strike it against a hard surface and state a judgment, hostile NPCs within Near range must test Resolve (DV 4) or hesitate for one exchange. Usable once per scene.

\section{Advanced Talents (6−−8 XP)}

≤galitem{Right of Acknowledgment (6 XP)} \textit{Requirement: You have witnessed an injustice involving kin or oath.} Once per arc, you may demand that an NPC speak one true answer to a question of your choice. The NPC may refuse, but refusal forfeits all legal protections (they become \textit{anima non grata} in Viterran territory). The truth need not be complete, but it must not be a lie.

≤galitem{Justicar's Immunity (8 XP)} You have been granted the Queen's own seal. You cannot be arrested or detained by any lesser authority in Viterra. However, abuse of this immunity (committing a crime in plain sight) will cause the seal to crack−−−lose this talent until you undergo a ritual of purification (a downtime project).

ørnament

\chapter{Patrons of the Viterran Courts}

\section{The Witness (Viterran Variant)}

Here, The Witness does not reveal hidden truths−−−it \textit{attests} to sworn testimony. Its shrines are courthouses; its rites are oaths spoken under a single lamp.

≤galitem{Low Rite −−− Testimony (4 XP)} When you repeat a statement you heard under oath, you may roll Presence + Investigation (DV 3). On a hit, you know whether the statement was false. (Mark +1 Obligation)

≤galitem{Standard Rite −−− Unbreakable Vow (8 XP)} You and a willing partner swear an oath before witnesses. Until the oath is fulfilled, both gain +1 die to any action that serves it. Breaking the oath causes 2 Harm (Spirit). (Mark +2 Obligation)

\section{Mykkiel (Judge of the Scar)}

Mykkiel's Viterran face is that of a circuit judge, riding from town to town, carrying a ledger of unpaid debts. His followers are Justicars who have lost their faith in mercy but not in order.

≤galitem{Low Rite −−− Summons (4 XP)} Name a person within Far range. They hear your voice as if you stood beside them. They are not compelled to obey, but they know the message is legal. (Mark +1 Fatigue)

≤galitem{Standard Rite −−− Attachment (8 XP)} Touch an object or person. For one day, any attempt to move them from their current location suffers +1 DV. (Mark +2 Obligation)

\section{The Queen (Mortal Patron)}

Queen Veralyn herself can serve as a Patron for those sworn directly to her service. Her Rites are not magical−−−they are legal and social permissions.

≤galitem{Low Rite −−− Royal Favor (4 XP)} Present a token with the Queen's seal. For one scene, any Viterran official treats you as if you were a royal messenger. (Mark +1 Exposure −− the Queen's enemies take note)

≤galitem{Standard Rite −−− Pardon (8 XP)} Forgive one crime committed by another person, as long as that crime did not harm the Queen's interests. The pardon is legally binding. (Mark +2 Obligation to the Queen; she may call on you later)

ørnament

\chapter{The Justicar's Judgment}

\begin{quotebox}
\textit{The tale of Baron Aldric the Turncloak, set down by Dusana of the Raven Road after a cup of small beer in a Viterran hedge−court. The ink is gall and the memory of a father's silence.}
\end{quotebox}

\vspace{2cm}
\begin{center}
\textit{A tale of Viterra, told by a clerk who sharpened the quills}
\end{center}
\vspace{1cm}

\epigraph{The law does not forget. Neither does the son who was not ransomed.}{−−− Queen's Justicar Elodie Venn, from the bench}

\section*{Summons – The Hedge−Court}

The evening mist clung to the hedge−court like a shy lover, curling around the low stone benches and the half−filled mugs of small beer that steamed faintly in the chill.  
Dusana set a cup down, the amber liquid catching the last glint of daylight, and laid before her a vellum sheet black−waxed with the crown of Viterra and the broken column of Vhasia.  

\bigskip
\textit{You are commanded to hear the case of Baron Aldric of the Scar…}  

\bigskip
Her hand did not tremble.

She remembered Sir Aldric of the White Moor−−−a laughing knight whose sword sang as he rode, whose grin could coax a smile from a storm. She had borne his shield at the Ford of Ashes, watched him barter a burnt tower for a handful of serfs, and seen him abandon Queen Veralyn's mother at the Battle of the Three Bridges, fling his banner to the Pretender's side, then swear fealty again when the Pretender fell.  

Now the queen's justice asked her to weigh the man who had once been her mentor against the scales of law and memory.  

≠wpage

\section*{The Justicar Arrives}

A pair of iron doors slammed open with a sound that rang like a distant clapper of a bell.  
Through the doorway stepped a figure that seemed to have been hewn from the very stone of the palace.  

Elodie, the Queen's Justicar, moved across the flagstones with the measured stride of a veteran knight.  Her cloak of deep violet bore the silver argent of the crown, the hem trimmed in the crimson of the queen's order.  A scar—white as bone—ran down the left side of her cheek, a souvenir from the Siege of Darwyn when a falling stone had caught her cheek as she shielded a comrade.  Her gauntleted left hand rested on the polished bronze gavel that perched on the bench, the metal catching the flickering torchlight.  

She was in her early forties, her hair pulled back into a tight braid that revealed a stern face framed by a crown of short, steel−gray hair.  In the years before she wore justice's seal, she had been squire to Sir Garen, a knight of the queen who had fought side−by−side with Aldric in the early wars of the Crown.  Garen's stories of Aldric's brilliance−−−of his youthful flash at the battle of Ember Hill−−−had filled the squeaking wooden walls of the keep where she learned swordplay.  She had never seen Aldric as a boy, only as the proud, reckless knight who'd once shared a campfire with her master.  When Garen fell at the Siege of Darwyn, she swore an oath to channel his memory of honor into the law, rising through the ranks of the Queen's Guard until the crown placed the Justicar's mantle upon her shoulders.

She paused at the center of the hall, the murmur of the gathered nobility falling to a hush.  The queen's herald stepped forward, his voice a thin ripple across the vaulted space.

\begin{description}
ℹtem[Herald:] ``By command of Her Majesty Queen Veralyn, the Court of the Hedge−Court is now in session. The Justiciar, Elodie of the Silver Gavel, shall preside.''
\end{description}

Elodie lifted the gavel, listening to its heavy ring echo off the stone.  In that moment the past and present pressed together: the knight she once admired, the baron now before her, the law she now embodied.  She knew the law could be a blunt instrument, but she also knew it could be sharpened by the will of the one who wielded it.

\begin{description}
ℹtem[Elodie:] ``Let the record note that the court is convened to hear the charges laid against Baron Aldric of the Scar, a man whose oath to the Crown has been called into question. The law shall be our guide, even when its reach feels insufficient.''
\end{description}

A murmur of assent rose from the benches.  The Justicar turned to the clerk who bore the heavy \textit{Codex of Precedents}, a tome as ancient as the kingdom itself.

≠wpage

\section*{Echoes of the Baron's Shadow}

The clerk placed the worn vellum upon the central table.  In the candle−lit margins were listed the accusations that had brought Aldric before the queen's justice.

\begin{description}
ℹtem[Clerk (reading):] ``Count One – Treason by repeated oath−breaking. Count Two – Unlawful levy of grain and levy of the Iron Tithe upon the peasants of the Burnt Tower. Count Three – Assault upon a servant, the loss of a hand, and the taking of a personal token without due process. Count Four – Refusal to ransom the firstborn son, Ser Roran.''
\end{description}

Elodie's gaze slipped briefly to the balcony where a solitary figure in grey observed the proceedings.  It was the wetnurse, Bryn, whose presence would soon become a stark illustration of the law's limits.

The Justicar's mind turned over the old maxim etched into the Codex: \textit{“The Crown may judge the noble, yet the lord's private jurisdiction over his own fief shall not be usurped without cause.”}  It was a reminder that some of Aldric's offences—particularly those he perpetrated on lands he claimed as his own—lay beyond the formal reach of the queen's courts.  

She inhaled, feeling the weight of the gavel in her hand, and spoke.

\begin{description}
ℹtem[Elodie:] ``The Court shall hear each charge in turn. The first three—treason, unlawful levy, and assault—are within our jurisdiction. The fourth, concerning the ransom of a son, falls into a grey area of custom and the Crown's lesser statutes. We shall treat it as a matter of moral reckoning, not strictly legal, and shall see what testimony the witnesses may provide.''
\end{description}

A thin, tense smile crossed the face of Aldric's counsel, but Elodie's eyes remained steady.  She knew the law could protect the powerful; she also knew that the law could be bent, if only a single voice dared to speak the truth.

≠wpage

\section*{Interlude – Eleanor, Daughter of the Accused}

The murmurs of the courtroom fell away for a heartbeat as the light from the high windows brushed the marble floor. In a quiet alcove near the rear of the hall, a figure turned the page of a thin vellum with a deliberate slowness that seemed to count the passing seconds more than the words.

Eleanor, the eldest child of Baron Aldric, sat beside a man whose coat bore the insignia of the House of Riven—her husband, the lord to whom she had been bound twenty−odd years earlier as a bargaining chip in Aldric’s ascent to power. The pairing was as ordinary in Viterra as a merchant’s contract; it required no testimony, no oath, only the quiet acknowledgment that alliances were forged on the anvil of ambition.

Eleanor’s once−sickly childhood had left her with a slight hollowness to her cheekbones, a softness that lingered like a whisper even as she had grown into a woman of striking beauty. Her hair, once the colour of sun−warmed wheat, now fell in loose, silver−streaked braids that framed a face etched with fatigue. Deep lines gathered at the corners of her eyes, not from age alone but from the weight of years spent watching a father’s rise and fall, from the endless counsel of courtiers, from the perpetual undercurrent of danger that lingered wherever Aldric’s name was spoken.

Across the aisle, three young men stood together—her half‑brothers, Aldric’s sons by a later marriage. They bore the same sharp jaw, the same cold confidence, the same smirk that had once adorned their father’s face before the wars wore it away. Eleanor did not look at them. She had not looked at them in years. To meet their eyes would be to see Aldric’s ghost already standing among the living, arrogant and unrepentant, and she had spent too long learning to look away.

Her husband, a tall, gaunt man in a dark, fur‑trimmed mantle, sat rigidly beside her, his hands clasped over a goblet of dark wine. He glanced intermittently at the proceedings, his gaze sharp, yet he offered no words—his role was that of a silent sentinel, a reminder that the Baron’s bloodline remained intertwined with the very nobility that now judged him. The faint scar on his left forearm, barely visible beneath his sleeve, caught the flicker of torchlight, a small testament to battles fought on behalf of his wife’s family and the price that had been paid for their union.

Elodie, seated a few rows forward, caught the tableau in a fleeting glance. The Justiciar’s eyes lingered on Eleanor’s weary posture, on the way the lady’s shoulders seemed to carry an invisible load heavier than any lace or jewel could conceal. In a moment’s pause, Elodie thought of the lawbooks stacked upon the bench—pages that spoke of treason, of levy, of the Right of Acknowledgment—yet ignored the quiet suffering that lived in the faces of those bound to the accused. The sight reminded her that the scales of justice, no matter how precisely calibrated, could not alone balance the personal toll exacted by the very decisions the law sought to adjudicate.

A soft sigh escaped Eleanor’s lips, barely audible over the rustle of parchment. She pressed a thin, tarnished ring—a token given to her on her wedding day—against the vellum, as if trying to anchor herself to a memory that might still hold a fragment of hope. The lord husband’s fingers tightened around the goblet, his knuckles whitening, a silent affirmation that he too felt the strain of a union born of politics rather than affection.

The courtroom resumed its rhythm, the witnesses took their places, and the judges returned to their deliberations. Eleanor remained seated, a stoic silhouette amidst the storm of accusations, her presence a quiet counterpoint to the clangor of gavel and oath—a living reminder that the law’s reach extended far beyond the statutes written in ink, touching the lives of those who, though they did not speak, bore its consequences with a dignity that spoke louder than any testimony.

≠wpage

≠wpage
\section*{Witness: Bryn, the Wetnurse (The Hooked Testimony)}
\smallskip
\textbf{Call the witness: Bryn, former wetnurse to Lady Margot, first wife of the accused.}

The gallery parted slowly. In the wavering torchlight a woman in a threadbare grey robe leaned on a knotted staff and moved toward the bar. As she passed, an inexplicable chill brushed the faces of those nearest; candle flames dipped, then steadied. Her left hand ended not in flesh but in a brass hook—rough, forged from a farmer’s rasp, still scarred by the fire that had shaped it. The hook held a faint, cold iridescence, as if a sliver of otherworldly light ran beneath the metal.

Her face was a map of old grief, skin drawn tight over sharp bones, and she seemed older than the stone of the courthouse itself. Yet when her eyes found Aldric they burned with a light that had not dimmed in forty years. Behind her, the polished marble reflected only the torches; she cast no shadow.

\begin{description}
ℹtem[Elodie:] ``State your name and relation to the accused.’’

ℹtem[Bryn (voice thin as winter ice, yet clear):] ``I am Bryn. I served as wetnurse to his first wife, Lady Margot, and to their daughter, Eleanor. I was midwife when she died.’’

ℹtem[Aldric’s counsel (rising, a sneer in his tone):] ``A servant’s testimony from forty years past? The woman is ancient, her memory clouded by age and spite. Moreover, the Baronet held pit and gallows on his own land. Any discipline he administered there is not subject to this court’s review.’’

ℹtem[Elodie (after a long pause, her gavel resting on the bench):] ``The court will hear the witness, though the baron’s right to enforce law on his holding is noted. Proceed.’’
\end{description}

Bryn’s good hand reached into her grey robe and drew out a silver locket, tarnished but still whole. As she lifted it, the air around her seemed to grow colder, and the torchlight dimmed a shade.

\begin{quote}
``Margot gave me this on her deathbed. ‘For the child,’ she said. ‘The one who still breathes.’ Aldric never asked for it. He never visited her grave. I left his service in disgust. Three days later his men found me at a waystation. He accused me of thieving – of stealing this locket, which he had never once mentioned while his wife lived. He did not ask. He just swung. He laughed and said the hand was enough; I could keep the ‘trinket.’’’
\end{quote}

She raised her brass hook toward the bench. The light caught its rough surface, and for an instant the hook glowed with a cold, faint fire that did not belong to any ordinary metal.

\begin{quote}
``This was forged from a farmer’s rasp, at a forge that no longer stands. The hand that fed his children, that held his dying wife’s head – he took it for a trinket he had forgotten existed. I left the Cord behind. But I do not forget. Neither does the law, if it has the courage to remember.’’
\end{quote}

A heavier murmur passed through the gallery. Some nodded; others looked away. Aldric’s counsel objected again, citing the baron’s right of pit and gallows. Elodie, after a whispered consultation with the clerk of the \textit{Codex of Precedents}, ruled the testimony admissible only as character evidence; the baron’s disciplinary authority on his own land could not be challenged by the Queen’s court. The substance of the accusation—the loss of Bryn’s hand—was dismissed as beyond the court’s jurisdiction.

But the locket remained on the evidence table, its tarnish catching a stray spark that seemed to linger after the torch's flame had faded. Bryn did not weep. She lowered her hook and limped back into the gallery, her eyes never leaving Aldric's face. He did not look at her. He could not. Yet his lips twitched into a thin, almost amused smile, as if he sensed something unspoken in the cold that followed her out of the hall. For a hair's breadth, Elodie thought she saw a young girl take the old woman's hand as she disappeared into the darkness.
≠wpage

\section*{Witness: Brother Tomas, Keeper of Saint Mirelle's Chapel (The Chapel−Burner)}

\textbf{Call the witness: Brother Tomas, keeper of Saint Mirelle's Chapel.}

The light was already slipping toward dusk when Tomas shuffled forward, his robe frayed at the hem, eyes haunted by the memory of flame. He took the oath, his voice a thin reed trembling in the gathering gloom.

\begin{description}
ℹtem[Elodie (softly):] ``Brother Tomas, tell the court what you saw on the night the chapel burned.''
ℹtem[Tomas (voice trembling):] ``My lord, I was at vespers when the torches came—smoke already curled from the nave like a warning. Men in the Baron's colours shattered the door, shouting that rebels hid within. They thrust flaming brands into the straw, and the fire rose like a hungry beast, licking the beams and devouring the rafters. I fought my way through the heat, pulling two children from the pews before the ceiling gave way; twelve souls were lost that night, sir.''
ℹtem[Aldric's counsel (leaning forward, voice oily):] ``And you, brother, were you not once a tenant of the Baron's lands? Might your testimony be coloured by a grudge?''
ℹtem[Tomas (flushing, hands clasping the rosary tighter):] ``I bear no grudge, only the memory of the innocent who perished. I speak for them, not for myself.''
\end{description}

The court murmured, the weight of those twelve lives settling like ash on the shoulders of those who heard, the fading light casting long, somber shadows across the stone floor.

≠wpage

\section*{Witness: Mirella, Widow of the Mill of the Burnt Tower (The Iron Tithe)}

\textbf{Call the witness: Mirella, widow of the miller of the burnt tower.}

Dawn had just broken, a thin veil of mist clinging to the fields as Mirella stepped forward, her hands calloused from years of grinding grain, her eyes hard as flint against the pale light.

\begin{description}
ℹtem[Elodie:] ``Mirella, describe the levy the Baron imposed on your household.''
ℹtem[Mirella (steady, voice low but resonant):] ``He arrived with his men at first light, declaring that protection required half our harvest. We gave it, trusting his word, believing the promise of safety. Then he carted the grain to market, selling it at full price while we watched our stores dwindle. My children grew thin; my husband fell ill from the strain, his cough echoing in the empty barn.''
ℹtem[Aldric's counsel (smirking, voice smooth as polished stone):] ``Was it not the Baron's right to demand tribute in times of war? Did you not benefit from the peace his forces brought?''
ℹtem[Mirella (voice rising, a flash of anger breaking through the dawn calm):] ``Peace bought with our blood is no peace at all. He took our sustenance to fund a private militia that terrorised the very folk he swore to protect. Each sack of grain felt like a blade against our throats.''
\end{description}

A ripple of disapproval passed through the nobles; the crisp morning air seemed to carry the metallic tang of greed, sharpening the sense that justice, though slow, was finally stirring.

≠wpage

\section*{Witness: Ser Roran, the Unransom'd Son (The Son Who Was Not Ransomed)}

\textbf{In the year of the Broken Crown, after the Battle of the Red Ditch, your son—Ser Roran, then seventeen—was captured by the forces of the Duke of Adria. The customary ransom for a knight's eldest son was three hundred silver marks. The Duke's herald offered you a discount of fifty marks if you paid within the week.}

Aldric's smile thinned to a razor's edge.

\begin{quote}
``I recall.''
\end{quote}

\begin{quote}
``You said, and I quote the herald's sworn testimony: ‘I have more sons and can make further. The boy was weak. Let the Duke keep him.’''
\end{quote}

A rustle stirred the gallery. In the shadows sat a lean, grey−faced man—Ser Roran, now thirty−seven, his left hand sheathed in a silvered gauntlet, his right resting on the pommel of a sword that had seen twenty campaigns.

Elodie continued, her voice low but resonant, each syllable a thread pulling the tapestry tighter.

\begin{quote}
``The Duke did not keep him. He threw the boy into a debtor's cell, expecting you to relent. You did not. For three years Roran rotted. Then a stranger—once a hedge−knight under your banner—paid the ransom with every coin he had, died of fever the following winter, and gave Roran his name and his sword.''
\end{quote}

She paused, letting the image of the boy's lonely cell settle like frost on stone.

\begin{quote}
``He did not return to you. He fought in the eastern wars. He won the Queen's own pennon at the Siege of the Horn. He knighted himself with a broken blade and the blessing of a battlefield priest. He has not spoken your name in seventeen years.''
\end{quote}

Aldric's jaw tightened, the muscle jumping like a trapped hare.

\begin{quote}
``The boy survived. Stronger for it.''
\end{quote}

Ser Roran rose. The hall fell silent as he walked to the bar, placed his gauntleted hand on the rail, and spoke, his voice barely above a whisper yet carrying the weight of years.

\begin{quote}
``Father, do you know what I remember most? Not the cell. Not the chains. The moment I heard you say it. I was hiding behind the herald's tent. I heard you laugh.''
\end{quote}

≠wpage

\section*{The Justicar's Verdict (The Judgment)}

The law of Viterra stood ancient and immutable: treason meant death; oath−breaking stripped lands and titles; refusing a son's ransom was not a crime—only cruelty.  

Elodie spent three hours upon the bench, the massive \textit{Codex of Precedents} hauled by two clerks like a tombstone. She asked for water, for silence, for the space to let the truth settle like dust after a storm.

When she rose, her voice was level, each word a stone placed carefully upon the scales of justice.

\begin{quote}
``Baron Aldric of the Scar, on the first count—treason by repeated oath−breaking—I find you guilty. The sentence is forfeiture of all lands, titles, and pensions. You will be stripped of your banner and your ring. You will live as a hedge−knight if any will take you.''
\end{quote}

Aldric's face remained unmoved; he had expected this.

\begin{quote}
``On the second count—refusal of ransom and abandonment of kin—the law is silent. Viterra does not compel a father to love. But the law does possess a clause for debts of the heart: the \textit{Right of Acknowledgment}. No one has invoked it in two hundred years.''
\end{quote}

She turned her gaze to Ser Roran.

\begin{quote}
``Ser Roran, son of Aldric, you hold the right to demand that your father speak one truth aloud, in this court, under witness. Should he refuse, he is declared \textit{anima non grata}—his soul outlawed from all Vritann burial grounds. His body will be interred at a crossroads with a stake through his chest. The law permits it; the Church condemns it. Yet the law is the law.''
\end{quote}

Aldric laughed—a short, bitter bark.

\begin{quote}
``A truth? I have spoken a thousand truths. I have said a man must cut what he cannot carry. I have said sons are arrows—you loose them and you do not weep. That is my truth.''
\end{quote}

Ser Roran stepped forward, his voice soft as falling silk yet unyielding.

\begin{quote}
``No. Your truth is that you were afraid. Not of the Duke. Of me. I was stronger than you at seventeen. You saw it. You could not bear to pay my ransom because then you would owe me, and you have never owed anyone in your life.''
\end{quote}

Aldric opened his mouth, closed it. The hall held its breath, the silence thick as the velvet drapes.

\begin{quote}
``I am not asking you to love me,'' Roran said. ``I am asking you to say it. Just once. Say: ‘I was afraid.’''
\end{quote}

Aldric stared at his son. The mask of the turncloak—the smile, the shrug, the armor of indifference—cracked like thin ice under a spring thaw.

\begin{quote}
``You killed my wife,'' he growled. ``You wear her face.''
\end{quote}

Roran stared at him, hard, waiting.  \textit{Liar.} The accusation was written in his eyes.

\begin{quote}
    ``I was afraid,'' he finally whispered.
\end{quote}

The sound was so small the clerks leaned forward, straining to catch it.

\begin{quote}
``Of what?'' Elodie prompted.
\end{quote}

\begin{quote}
``Of him. Of what he would become. He would have been better than me. And I… I could not stand to watch.''
\end{quote}

Ser Roran turned away, his steps silent upon the stone, the silver gauntlet clicking a faint, steady rhythm. He did not weep; the tears had long since dried into salt upon his cheeks.

Elodie lifted her gavel.

\begin{quote}
``The Right of Acknowledgment is satisfied. The court is adjourned. Baron Aldric—former Baron Aldric—you are free to go. But I advise you to leave Viterra by sunset. The road is long, and the winter is coming.''
\end{quote}

≠wpage

\section*{Epilogue – Aftermath}

Three winters later Aldric was found dead in a ditch outside Thepyrgos, near the Black River, reeking of wine. His lay body unclaimed, circled in salt—he had outlived his own burial rights. Scratched into his chest were the words \textit{anima non grata} and a strange, circular wound near his heart.

Ser Roran entered the Queen's household as a knight, wearing the silvered gauntlet over his left hand—not to conceal the scar, but to remind himself that a stranger's coin had bought his life. He never took a ransom. Once a captured lord offered him five hundred marks for his freedom.

\begin{quote}
``Pay it to the hedge−knight's widow,'' he said. ``I have no need for gold. I had a father once. That was expensive enough.''
\end{quote}

When Queen Veralyn learned of the verdict, she sat in silence for a long breath, then summoned the \textit{Codex of Precedents} and, in her own hand, added a marginal note that would become law:

\begin{quote}
``The Right of Acknowledgment may be invoked by any child who was not ransomed, not visited, not named in a will. The truth, once spoken, is payment in full. No further interest shall accrue.''
\end{quote}

The scribes copied the clause into every new edition. It remains law to this day.

\noindent\textit{The ink is dry. The gavel is still.}

ørnament

\chapter{Adventure Seeds}

\begin{enumerate}
    ℹtem \textbf{The Unfinished Testament.} A dying noble calls for a Justicar to witness his will. The will disinherits his eldest son in favor of a bastard daughter who was never ransomed from captivity. The son hires thugs to intercept the Justicar. The party must deliver her to the bedside before the old man dies.
    
    ℹtem \textbf{Hedge−Court at Dusk.} A village accuses a turncloak knight of stealing a boundary stone. The knight claims the stone was always on his land. The party must survey the hedge, interview the oldest residents, and render a verdict that will not start a blood−feud.
    
    ℹtem \textbf{The Turncloak's Daughter.} A Vhasian baron's daughter, now living in Viterra, petitions for the Right of Acknowledgment against her father's estate. The father is dead, but his executors refuse to release her inheritance. The party must find a witness who heard the father speak a damning truth−−−or forge a confession that the hedge−court will accept.
    
    ℹtem \textbf{The Codex Theft.} A page from the \textit{Codex of Precedents} has been stolen−−−a page that would decide a pending lawsuit between two powerful houses. The party must track the thief through the underworld of legal scribes and forgers before the case is dismissed for lack of precedent.
    
    ℹtem \textbf{The Queen's Suitors.} Queen Veralyn has announced she will consider marriage offers. Every noble house in Viterra sends a champion to court her. The party is hired by a minor house to sabotage the other suitors−−−legally, of course. No poison. No daggers. Just clever readings of dowry law, ambush by precedent, and a well−timed Right of Acknowledgment.
\end{enumerate}

ørnament

\chapter{GM Toolkit: Timers and Fronts}

\section{Sample Timers}

\begin{longtable}{@{}p{2.5cm}p{4cm}p{5cm}@{}}
⊤rule
\textbf{Clock} & \textbf{Purpose} & \textbf{Advance / Retreat} \\
∣rule
\endhead
\textbf{The Queen's Patience [6]} & How long the crown tolerates noble infighting & Advance when parties flout Justicar rulings; retreat when a major case is settled cleanly \\
\addlinespace
\textbf{Hedge−Law Fracture [4]} & Local customs breaking down under pressure & Advance when a hedge−court's verdict is ignored; retreat when a respected elder restores peace \\
\addlinespace
\textbf{Turncloak Conspiracy [8]} & A plot among pardoned nobles to reclaim lost lands & Advance when the party ignores rumors or protects a turncloak; retreat when a conspirator is exposed \\
\addlinespace
\textbf{Codex Decay [6]} & The \textit{Codex} losing authority as precedents contradict & Advance when contradictory rulings are issued; retreat when a Justicar reconciles two precedents \\
⊥tomrule
\end{longtable}

\section{Sample Front: The Succession Crisis}

\textbf{Impulse:} The Queen has no heir. Her marriage negotiations have stalled. The turncloak nobility smell blood.

\textbf{Grim Portents:}
\begin{itemize}
    ℹtem A hedge−court declares a local lord has the "ancient right" to name the next Queen.
    ℹtem The Parlement refuses to fund the Justicars, leaving a key case unheard.
    ℹtem A turncloak knight invokes the Right of Acknowledgment against the Queen herself, demanding she name her father's true heir.
\end{itemize}

\textbf{Timers:}
\begin{itemize}
    ℹtem \textit{Throne Vacancy [8]} −− When full, a claimant marches on the White Hall.
    ℹtem \textit{Justicar Defection [4]} −− When full, one of the twelve Justicars declares for a rival.
    ℹtem \textit{People's Petition [6]} −− When full, a hedge−court of commoners issues a binding judgment on the succession, overriding noble law.
\end{itemize}

ørnament

\chapter{Colophon}

This expansion was written in a Viterran hedge−court, on paper that had been used to wrap cheese. The ink is gall and the ash of a burned writ. The binding is leather from a debtor's cow.

The Queen's Justicar Cass Venn died in her sleep seven years after the trial of Baron Aldric. Her gavel is kept in the White Hall, next to the \textit{Codex of Precedents}. Children who visit the Hall are told to touch it for luck. They do.

Ser Roran served the Queen for thirty years. He never married. He never took a squire. He wore the silver gauntlet until the day he died, and it was buried with him. The hedge−knight who ransomed him has no grave−−−his body was never found. But Roran paid his widow every year on the anniversary of the ransom, in full, with interest.

The road to Viterra is straight, the hedges are trimmed, and the law is patient.

\begin{center}
\vspace{2em}
\textit{The ink is dry. The gavel is still.}
\end{center}

\end{document}